\documentclass[12pt]{article}
\usepackage[utf8]{inputenc}
\usepackage[T1]{fontenc}
\usepackage{hyperref}
\title{Tzotzil (Chenal\'o) Language Resources for K-12 Curriculum Development}
\author{Compiled by Claude Sonnet 5 (claude-sonnet-5) for Andr\'as Kornai}
\date{}
\begin{document}
\maketitle
\section{Introduction}
Tzotzil, Chenal\'o (ISO 639-3 tze, Glottocode tzot1263) is a Mayan variety of Chiapas, Mexico. The identification line supplied with the task gave the family as ``Bookkeeping''. That appears to be a metadata error, since Tzotzil is Mayan. Search turned up little that is specific to the Chenal\'o variety. Most substantial published work covers the Zinacant\'an and San Bartolom\'e dialects.
\section{Primary Educational Resources}
No Chenal\'o primer, textbook or school programme turned up in search.
\section{Linguistic and Reference Materials}
\begin{itemize}
\item Great Tzotzil Dictionary of San Lorenzo Zinacant\'an (Smithsonian). It covers a neighbouring dialect, not Chenal\'o. \url{https://repository.si.edu/handle/10088/1360}
\item Smithsonian repository record for a companion Zinacant\'an dictionary volume. \url{https://repository.si.edu/handle/10088/1367}
\item A Descriptive Grammar of the Tzotzil Language as Spoken in San Bartolom\'e de los Llanos, Chiapas (Sarles, 1966 dissertation). It covers a different dialect. \url{https://www.proquest.com/docview/302240547}
\item General overview of the Tzotzil language and its dialects. \url{https://en.wikipedia.org/wiki/Tzotzil_language}
\end{itemize}
\section{Religious and Community Resources}
Global Recordings lists evangelistic and Bible-teaching audio in Tzotzil: Chenalho. \url{https://www.globalrecordings.net/en/language/2681} The search summary also pointed to Scripture Earth scripture for Chenal\'o, but I did not see its URL.
\section{Digital and Audio Resources}
The Living Dictionaries Tsotsil entry I saw is a short place-name entry, not a general dictionary of Chenal\'o. \url{https://livingdictionaries.app/batsi-kop-tsotsil-tsot/entry/VyPuoeNawoWHbcFjzBsj}
\section{Curriculum Development Context}
Little exists that is specific to Chenal\'o. A curriculum would need to adapt Zinacant\'an-based materials, or create new materials with Chenal\'o speakers. No information on bilingual schooling for this variety turned up in search.
\end{document}